\section{Introduction}\label{sec:intro}

An LLM agent with tools operates as a confused deputy in the textbook sense:
it acts with the authority of its principal while consuming data and
instructions across shared channels. Whenever an agent processes sensitive
context alongside unvetted tool outputs, external inputs, or untrusted data,
structural risks inevitably arise. Unauthorized data exfiltration, unintended
tool calls, and policy violations can result not only from adversarial prompt
injection---where attacker-controlled text hijacks control flow---but also from
model hallucinations, ambiguous instructions, or miscalibrated agent reasoning.
In all these cases, exposing a capable agent to private data, external sinks,
and unvetted content turns read access into a vector for exfiltration and policy
breaches.

Many deployed defenses cluster at two poles. At one pole, imperative
guardrails: prompt-injection classifiers, regex filters, and output hygiene.
These judge \emph{content}, so they are bypassed by paraphrase, offer no
guarantees against non-adversarial agent errors, and treat prompt injection as an
isolated bug rather than a symptom of mixed-confidentiality processing. At the
other pole, hard restriction: allowlists, human approval on every call, and
no-tools-after-untrusted-read policies. These judge \emph{structure} but do so
statically and coarsely, taxing every benign action and destroying agent
utility.

Agent information-flow control (IFC) offers a third approach. CaMeL separates
control and data flow~\cite{camel2025}; Fides propagates confidentiality and
integrity labels through messages, actions, arguments, and results, checking
consequential calls against deterministic policies~\cite{fides}. Fides
also identifies the utility problem central to this paper: allowing a
restrictive result to enter planner context can permanently disable useful subsequent
calls. Fides addresses this by selectively storing such results in labeled variables
and revealing constrained derivatives through a quarantined LLM. \appa{} builds
directly on this premise. Rather than claiming deterministic IFC enforcement or
selective variable hiding as novelties, \appa{} asks whether a restrictive read
can be explored without tainting the long-lived planner trajectory at all.

Formally, \appa{} maintains execution state using two core abstractions.
First, a \textbf{label}---defined as an audience (a reader set) $\times$ trust (a rank
in a finite chain)---accompanies data and is verified at policy checks. Tool contracts
declare label \emph{actions} (restrictive reads intersect reader sets and take minimum
trust levels), and the overall trajectory label is the composition of actions from all
dispatched calls. Second, an append-only \textbf{event log} records execution history
without requiring separate approval. The log captures world events (such as external egress
or state mutations) upon dispatch, as well as governance events (rulings, boundary markers,
sanitizer invocations, and label casts). Every policy decision reduces to a label comparison,
a log predicate, or a declared per-call attention requirement; imperative logic---such as approval
workflows, sanitizers, and ACL resolution---resides in registered external components.
Because label actions only restrict access, label updates are monotone and commutative,
ensuring predictable, deterministic state aggregation.

The central mechanism of \appa{} is \textbf{branching as taint prevention}. In classical
IFC, read operations incur no immediate restriction, and enforcement occurs when sensitive
data reaches an external sink. \appa{} instead intercepts a prospective label descent prior to
the read operation and can fork a child trajectory. The child inherits the parent's label to
prevent policy evasion, and only the child receives the restrictive result. The child can inspect
the result and execute further tool calls under its descended label, while the parent trajectory
remains unaffected at its pre-fork label. At exit, the branch is either abandoned or returns a single
explicitly labeled value. Returning raw data propagates restrictions back to the parent via label
intersection; preserving the parent's broader access therefore requires a registered sanitizer or a
quarantine-exit attestation that bounds the derivative value. \appa{} never removes taint from an
already exposed trajectory; instead, it prevents exposure by confining restrictive descents to disposable
child branches.

This design transforms data acquisition into an explicit policy event. \appa{} orders label states
by restrictiveness and soft-blocks any tool call whose prospective state strictly descends in permission.
The resulting execution options present clear choices: accept the label narrowing in the current
trajectory, fork and confine the observation within a child branch, apply a registered sanitizer, or
decline the read. Because non-elevating steps only descend in permission and all remedies are statically
registered, the candidate plan space is finite and fully enumerable. Consequently, an empty candidate plan
set serves as a formal proof that no registered remedy exists for a given action, rather than an incomplete
search result.

Two structural rules ensure that branching preserves security guarantees and execution integrity. First, a child trajectory inherits its parent's label at fork time and any merged result updates the parent label via set intersection and minimum trust rank, preventing unvalidated data elevation. Second, all branches append real-world side-effects (such as external egress) to a single, globally shared event log in real time; even if a child branch is later abandoned, its real-world side-effects remain recorded. Thus, branching isolates trajectory taint without concealing executed actions.

Atomic, call-scoped judgment supports this architecture when discretionary release is required.
The engine renders the pending tool call with its fully resolved arguments and provenance (rather than
relying on agent paraphrases), submits the exact call to an authorizing policy component, and dispatches
it in a single indivisible step upon approval. Because approval and dispatch coincide, no reusable grant
tokens are created, preventing replay or misbinding attacks. Combined with the soundness of the candidate
plan set---where every proposed plan is individually safe regardless of which option a potentially untrusted
agent selects---this yields a robust declassification discipline for agentic workflows~\cite{zdancewic2001robust}.

Contributions, in order of importance:
\begin{enumerate}
  \item \textbf{Branching as taint prevention:} fork-before-read semantics
    that confine restrictive observations to a child, preserve the parent
    before merge, permit only policy-bounded exits, and retain honest world
    history in a shared log (\S\ref{sec:branching}).
  \item A prospective acquisition check that exposes acceptance, branching,
    sanitization, and refusal as sound plans before restrictive data enters a
    trajectory, with plans finitely enumerable from the registry
    (\S\ref{sec:enforcement}).
  \item A two-monoid model supporting the construction---label actions
    checked, events recorded---with policy as pure data and proofs of monotone
    descent, parent preservation, and merge confinement (\S\ref{sec:model}).
  \item Atomic, call-scoped rulings over engine-rendered requests, used for
    bounded branch exits and other exceptional dispatches without ambient
    elevation or label widening (\S\ref{sec:rulings}).
  \item An open-source implementation---a pure decision core under a
    state-owning shell---evaluated on
    AgentDojo~\cite{debenedettiAgentDojoDynamicEnvironment2024} and compared
    directly with Fides (\S\ref{sec:implementation}--\S\ref{sec:evaluation}).
    \todo{soften/sharpen once numbers exist.}
\end{enumerate}
